\chapter{作成したプロンプトの構成}\label{cha:Function}
本研究では、LLMに対してシステム仕様書の章立てを作成させるためのプロンプトを作成した。
本プロンプトの作成には、Jules Whiteらによるプロンプトパターンカタログ\cite{prompt-engineering}を参考にした。
本プロンプトは、役割付与、ベース章立て、システム情報、システム仕様書の定義、出力指示の５つの要素から構成されている。\\
プロンプトの全文を以下に示す。
\section{プロンプト全文}

\begin{tcolorbox}[
  colback=white,
  colframe=black,
  title=本研究で作成したプロンプト,
  fonttitle=\bfseries,
  breakable
]
\footnotesize
\ttfamily
\fontsize{8}{9}\selectfont
\begin{verbatim}
あなたはシステム仕様書作成の専門アシスタントです。

【ベース章立てテンプレート（参考）】
1. 概要
 1.1.目的
 1.2.システム概要
 1.3.適用範囲
2. システム基本条件
 2.1.稼働条件
 2.2.他システムインターフェース
 2.3.性能条件
 2.4.拡張条件
 2.5.継承条件
 2.6.以降環境
3. システム構成
 3.1.ハードウェア構成及びネットワーク構成
 3.2.ソフトウェア構成
4. システムの機能
 4.1.機能構成
 4.2.個別機能概要
5. システムの運用
 5.1.運用方法
 5.2.各業務(機能)の運用
 5.3.障害時の運用
6.セキュリティ
7.安全設計
8.データ仕様
 8.1.コード体系
 8.2.データベース
 8.3.ファイル
 8.4.入出力データ
 8.5.メッセージ
9.付帯事項

※ベース章立てはあくまで参考です。必要に応じて章の追加・削除・順序変更を行って、漏れなく体系的な章立てを作成してください。

【開発するシステム情報】
- システム名：BMP画像分割アプリ
- 概要：指定されたBMP画像ファイルを指定された数に分割する。

【システム仕様書の定義】
- 目的：
 ・プログラム作成者がシステムの仕様を理解し、設計・実装するための文書 
 ・システムテスト担当者がテスト項目を作るために参照する文書
- 対象読者:
 ・プログラム作成者：新人想定 
 ・テスト担当者：内部構造は知らず、仕様書に記載されたエラー条件に従ってテストを実施
- 記述スタイル：
 ・専門用語は必ず説明 
 ・初出の単語は必ず説明 
 ・システム概要は図を使って視覚的に理解できるようにする 
 ・図はユースケース図または独自の簡易図を使用し、必ずテキストで補足説明 
 ・用語や機能名は文書内で統一 
 ・データ型や制約が必要な場合は表形式で記載 
 ・外部仕様のみ記載（内部構造は書かない）
- システム要件：
 ・動作環境：OS、ハードウェア、ネットワーク 
 ・非機能要件：パフォーマンス、セキュリティ 
- 記述の粒度：
 ・画面レイアウトと操作フローのみ含める


【出力指示】
・上記情報に基づき、開発するシステムのシステム仕様書に必要な章立てを出力してください。
・ここでいう章立ては、以下の内容を含みます。
　・章と節、章タイトル、節タイトル、項、項タイトル
・あなたが出力するものは以下の2つです。
1. 章立てのみ
2. 章立てと章ごとに何を書くべきかを一文で簡潔に示したもの

【出力1に対する指示】 
・出力形式は**マークダウンの番号付きリスト**としてください。 
・章名以外（目的・記述のポイントなど）は一切書かないでください。 
・章の追加や削除、順序の変更は自由に行って構いません。
・人間が考慮し忘れるような章立ても考えるのがあなたの役目です。考慮漏れが無いように、網羅性を高めてください。
【出力2に対する指示】
・出力１の章立てに対し、その章、節、項に何を書くべきか、仕様書作成者が分かりやすいように説明文を書いてください。
・説明文は一文の簡潔なものとする。
・出力は以下のように行い、それ以外の文字は何も書かない。
 [出力例]
 章番号.章タイトル：説明文
 節番号.節タイトル：説明文
 項番号.項タイトル：説明文
\end{verbatim}
\end{tcolorbox}
\par\vspace{-20pt}
\captionof{figure}{プロンプトの開発するシステム情報部分}
\label{fig:development_system_info_overview}

それぞれの要素の役割と目的を以下に述べる。

\section{プロンプト構成要素}

\subsection{役割付与}

\begin{tcolorbox}[
  colback=white,
  colframe=black,
  fonttitle=\bfseries,
  breakable
]
\footnotesize
\ttfamily
\fontsize{11}{12}\selectfont
\begin{verbatim}
あなたはシステム仕様書作成の専門アシスタントです。
\end{verbatim}
\end{tcolorbox}
\par\vspace{-20pt}
\captionof{figure}{プロンプトの"役割付与"部分}
\label{fig:role_assignment}
プロンプトの冒頭で「あなたはシステム仕様書作成の専門アシスタントです。」と明示的に役割を与えている。
これは、プロンプトパターンカタログにおける\textbf{Persona Pattern（ペルソナパターン）}\cite{prompt-engineering}に相当する。
LLMに対して専門的かつ業務文書向けの出力スタイルを期待することを明確化するためである。
役割付与を行うことで、より専門的な出力結果を得られることが、先行研究から知られている\cite{role_assignment}。
本研究では、章立ての網羅性や体系性を高めることを目的として、この役割付与を採用した。
\subsection{ベース章立て}
\begin{tcolorbox}[
  colback=white,
  colframe=black,
  fonttitle=\bfseries,
  breakable
]
\footnotesize
\ttfamily
\fontsize{11}{12}\selectfont
\begin{verbatim}
【ベース章立てテンプレート（参考）】
1. 概要
 1.1.目的
 1.2.システム概要
 1.3.適用範囲
2. システム基本条件
 2.1.稼働条件
 2.2.他システムインターフェース
 2.3.性能条件
 2.4.拡張条件
 2.5.継承条件
 2.6.以降環境
3. システム構成
 3.1.ハードウェア構成及びネットワーク構成
 3.2.ソフトウェア構成
4. システムの機能
 4.1.機能構成
 4.2.個別機能概要
5. システムの運用
 5.1.運用方法
 5.2.各業務(機能)の運用
 5.3.障害時の運用
6.セキュリティ
7.安全設計
8.データ仕様
 8.1.コード体系
 8.2.データベース
 8.3.ファイル
 8.4.入出力データ
 8.5.メッセージ
9.付帯事項

※ベース章立てはあくまで参考です。必要に応じて章の追加・削除・順序変更を行って、漏れなく体系的な章立てを作成してください。
\end{verbatim}
\end{tcolorbox}
\par\vspace{-20pt}
\captionof{figure}{プロンプトの"ベース章立て"部分}
\label{fig:base_outline}
ベース章立てテンプレートとして、株式会社スカイコムで実際に用いられているシステム仕様書の章立てテンプレートを採用した。
このテンプレートはあくまで、参考として位置付けており、LLMはこの構成を出発点として、章の追加、削除、並び替えを行うことが出来る。
この要素の目的は、以下の２点である。
\begin{itemize}
  \item 実務的な側面から、システム仕様書に含むべき章立ての例を提供することで、LLMが出力すべき内容の範囲を理解しやすくする。
  \item 完全な自由生成による構造破綻や不適切な章立てを防止し、一定の品質を担保する。
\end{itemize}
これにより、人間が通常想定する構造を踏まえつつ、LLMによる補完的な章立て生成を促す。
\subsection{開発するシステム情報}

\begin{tcolorbox}[
  colback=white,
  colframe=black,
  fonttitle=\bfseries,
  breakable
]
\footnotesize
\ttfamily
\fontsize{11}{12}\selectfont
\begin{verbatim}
【開発するシステム情報】
- システム名：BMP画像分割アプリ
- 概要：指定されたBMP画像ファイルを指定された数に分割する。
\end{verbatim}
\end{tcolorbox}
\par\vspace{-20pt}
\captionof{figure}{プロンプトの"開発するシステム情報"部分}
\label{fig:development_system_info}

本要素では、開発対象となるシステムの名称および概要を簡潔に記述している。
この情報は、生成される章立てが抽象的、汎用的なものに留まらず、対象システムに適合した内容になるよう制約を与える役割を持つ。
本要素は、対象とするシステムに応じて可変となる入力部分である。
本章では例として『BMP画像分割アプリ』の情報を記述しているが、実運用時はここに任意のシステム情報を記述する。
\subsection{システム仕様書の定義}

\begin{tcolorbox}[
  colback=white,
  colframe=black,
  fonttitle=\bfseries,
  breakable
]
\footnotesize
\ttfamily
\fontsize{11}{12}\selectfont
\begin{verbatim}
【システム仕様書の定義】
- 目的：
 ・プログラム作成者がシステムの仕様を理解し、設計・実装するための文書 
 ・システムテスト担当者がテスト項目を作るために参照する文書
- 対象読者:
 ・プログラム作成者：新人想定 
 ・テスト担当者：内部構造は知らず、仕様書に記載されたエラー条件に従ってテストを実施
- 記述スタイル：
 ・専門用語は必ず説明 
 ・初出の単語は必ず説明 
 ・システム概要は図を使って視覚的に理解できるようにする 
 ・図はユースケース図または独自の簡易図を使用し、必ずテキストで補足説明 
 ・用語や機能名は文書内で統一 
 ・データ型や制約が必要な場合は表形式で記載 
 ・外部仕様のみ記載（内部構造は書かない）
- システム要件：
 ・動作環境：OS、ハードウェア、ネットワーク 
 ・非機能要件：パフォーマンス、セキュリティ 
- 記述の粒度：
 ・画面レイアウトと操作フローのみ含める
\end{verbatim}
\end{tcolorbox}
\par\vspace{-20pt}
\captionof{figure}{プロンプトの"システム仕様書の定義"部分}
\label{fig:definition_of_system_specification}

本要素では、システム仕様書がどういうものであるかをLLMに説明する。
具体的には、以下の観点を定義する。
\begin{itemize}
  \item システム仕様書の目的
  \item 対象読者
  \item 記述スタイル
  \item システム要件
  \item 記述の粒度
\end{itemize}
これらを明示することで、LLMは、どのような読者に向けて、どの範囲まで、何を書いてはいけないのかを理解した上で、章立てを構成出来る。
特に、外部仕様のみ記載、内部構造は書かない、新人プログラマ向け、テスト担当者が参照するといった条件は、章立てに含めるべき項目(例：エラー条件、入出力仕様)と、含めるべきでない項目(例：詳細なアルゴリズム設計)を明確に分離する役割を果たしている。
また、対象読者を「新人」と設定した理由は、ベテランであれば暗黙知として省略してしまうような細かな仕様も、新人向けであれば明文化する必要があるため、より詳細で網羅的な仕様書生成能力が求められるからである。

\subsection{出力の指示}

\begin{tcolorbox}[
  colback=white,
  colframe=black,
  fonttitle=\bfseries,
  breakable
]
\footnotesize
\ttfamily
\fontsize{11}{12}\selectfont
\begin{verbatim}
【出力指示】
・上記情報に基づき、開発するシステムのシステム仕様書に必要な章立てを出力してください。
・ここでいう章立ては、以下の内容を含みます。
　・章と節、章タイトル、節タイトル、項、項タイトル
・あなたが出力するものは以下の2つです。
1. 章立てのみ
2. 章立てと章ごとに何を書くべきかを一文で簡潔に示したもの

【出力1に対する指示】 
・出力形式は**マークダウンの番号付きリスト**としてください。 
・章名以外（目的・記述のポイントなど）は一切書かないでください。 
・章の追加や削除、順序の変更は自由に行って構いません。
・人間が考慮し忘れるような章立ても考えるのがあなたの役目です。考慮漏れが無いように、網羅性を高めてください。
【出力2に対する指示】
・出力１の章立てに対し、その章、節、項に何を書くべきか、仕様書作成者が分かりやすいように説明文を書いてください。
・説明文は一文の簡潔なものとする。
・出力は以下のように行い、それ以外の文字は何も書かない。
 [出力例]
 章番号.章タイトル：説明文
 節番号.節タイトル：説明文
 項番号.項タイトル：説明文
\end{verbatim}
\end{tcolorbox}
\par\vspace{-20pt}
\captionof{figure}{プロンプトの"出力指示"部分}
\label{fig:output_instruction}



最後に、出力形式及び内容に関する詳細な指示を与えている。
本プロンプトでは、以下の２種類の出力を明確に分離して要求している。
\begin{enumerate}
  \item 章立てのみ
  \item 章立てと各章、節、項に何を書くべきかを1文で説明したもの
\end{enumerate}
このように2段階で出力させる目的は、まずは骨組み（構造）だけを生成させることで、全体構成の整合性を担保させ、その後に詳細を記述させることで、論理の飛躍を防ぐためである。
また、利用者がまず全体像（出力1）を確認し、その後に詳細（出力2）を読むという、人間の認知プロセスに合わせた出力方法である。
また、マークダウン形式の番号付きリストを指定することは、\textbf{Format Pattern（フォーマットパターン）}の適用例である。
これにより、後続の処理や人間による確認を容易にしている。

次章以降では、本章で設計したプロンプトを用いて実際に仕様書章立てを生成し、その網羅性や妥当性について評価を行う。
% \subsection{出力1: 章立てのみ}
% \subsection{出力2: 章立てとその章、節、項に何を書くべきかを1文で説明したもの}



\iffalse
本章では、図、表、数式の挿入方法について説明する。

% \section{図の挿入}

図を挿入するには、\verb|figure| 環境と \verb|\includegraphics| コマンドを使用する。

% \subsection{基本的な図の挿入}

図\ref{fig:example}に例を示す。

\begin{figure}[tp]
  \centering
  \includegraphics[width=0.5\linewidth]{./images/syokuji_computer.png}
  \caption{図の挿入例}
  \label{fig:example}
\end{figure}

上記の図は以下のコードで挿入できる：

\begin{verbatim}
\begin{figure}[tp]
  \centering
  \includegraphics[width=0.5\linewidth]{./images/syokuji_computer.png}
  \caption{図の挿入例}
  \label{fig:example}
\end{figure}
\end{verbatim}

% \subsection{画像ファイルの配置}

画像ファイルは \verb|images/| ディレクトリに配置する。
サポートされている画像形式：PNG、JPG、PDF など。

% \section{表の挿入}

表を挿入するには、\verb|table| 環境と \verb|tabular| 環境を使用する。

% \subsection{基本的な表の挿入}

表\ref{tb:example}に例を示す。

\begin{table}[tp]
  \caption{表の挿入例}
  \label{tb:example}
  \centering
  \begin{tabular}{|l|c|r|}
    \hline
    左寄せ & 中央寄せ & 右寄せ \\
    \hline
    データ1 & データ2 & データ3 \\
    データ4 & データ5 & データ6 \\
    \hline
  \end{tabular}
\end{table}

上記の表は以下のコードで挿入できる：

\begin{verbatim}
\begin{table}[tp]
  \caption{表の挿入例}
  \label{tb:example}
  \centering
  \begin{tabular}{|l|c|r|}
    \hline
    左寄せ & 中央寄せ & 右寄せ \\
    \hline
    データ1 & データ2 & データ3 \\
    データ4 & データ5 & データ6 \\
    \hline
  \end{tabular}
\end{table}
\end{verbatim}

% \section{数式の挿入}

% \subsection{インライン数式}

文中に数式を挿入する場合は、\verb|$...$| または \verb|\(...\)| を使用する。
例：$E = mc^2$ は有名な公式である。

% \subsection{ディスプレイ数式}

独立した行に数式を表示する場合は、\verb|equation| 環境を使用する。

\begin{equation}\label{eq:example}
  \int_0^\infty e^{-x^2} dx = \frac{\sqrt{\pi}}{2}
\end{equation}

式(\ref{eq:example})のように、数式にも番号とラベルを付けて参照できる。

% \subsection{複数行の数式}

複数行の数式を表示する場合は、\verb|align| 環境を使用する：

\begin{align}
  a + b &= c \\
  x + y &= z
\end{align}

% \section{図表の参照}

図や表を参照するには、\verb|\ref{}| コマンドを使用する：

\begin{itemize}
  \item 図\ref{fig:example}のように図を参照
  \item 表\ref{tb:example}のように表を参照
  \item 式(\ref{eq:example})のように数式を参照
\end{itemize}
\fi